\section{Horizon Chartreuse}

\subsection*{Deskripsi Soal}

Di lembar catatan riset berikutnya, Marisa menemukan laporan mengenai fenomena
aneh yang menimpa langit malam. Seseorang telah mencuri bulan purnama yang asli
dan menggantinya dengan purnama palsu yang redup. Laporan yang memuat kronologi
hilangnya bagian dari langit tersebut kini tidak dapat dibaca karena seluruh
teksnya dienkripsi menggunakan \textbf{Vigen\`ere \textit{Cipher}}.

Manfaatkan metode \textbf{\textit{Kasiski Examination}} untuk mengidentifikasi
perulangan pola karakter pada \textit{ciphertext} dalam memperkirakan panjang
kuncinya. Setelah panjang kunci diperoleh, terapkan analisis frekuensi bahasa
Inggris pada tiap sub-teks untuk merekonstruksi kunci dan memulihkan seluruh
isi naskah.

\medskip
\noindent\textbf{Berkas lampiran:}
\begin{itemize}[nosep]
  \item \textit{Attachment}: \url{https://drive.google.com/file/d/1LyCID57V3ckzlvQ9Wnj-mGg_YMHoP50e/view?usp=sharing}
  \item \texttt{md5sum}: \texttt{6b933809bbedb3c5c0f1fd9e2b1076cd}
  \item \texttt{sha256sum}: \texttt{d6ecf3775c56593f833db5c62e1447f2668b41a310562e7a9ccd6a1486d5f641}
\end{itemize}

\subsection*{Berkas \textit{Ciphertext}}

\ciphertextfile{../../src/22/v.txt}

\subsection*{Langkah Penyelesaian}

\textbf{Vigen\`ere \textit{Cipher}} rentan terhadap analisis \textbf{\textit{Kasiski Examination}} jika kunci yang digunakan relatif pendek. Metode ini mengasumsikan bahwa pengulangan rangkaian karakter (\textit{substring}) pada \textit{ciphertext} terjadi ketika \textit{substring plaintext} yang sama dienkripsi menggunakan bagian kunci yang sama. 

Penyelesaian soal ini dibagi menjadi dua tahap: \textbf{identifikasi panjang kunci} dengan \textit{Kasiski Examination}, lalu \textbf{pencarian kunci} dengan \textit{Frequency Analysis} bahasa Inggris.

\subsubsection*{Tahap 1: \textit{Kasiski Examination}}

\begin{enumerate}
  \item Cari semua rangkaian karakter yang berulang pada \textit{ciphertext} dengan panjang minimal 3 karakter (\texttt{VigenereSolver.find\_repeated\_sequences} pada \texttt{src/22/solver.py}).
  \item Hitung jarak antarkemunculan dari masing-masing \textit{substring} berulang tersebut.
  \item Cari faktor dari setiap jarak tersebut dan hitung frekuensi kemunculan setiap faktor \\
  (\texttt{VigenereSolver.get\_kasiski\_factors}). Faktor yang paling sering muncul berpeluang besar merupakan panjang kunci atau kelipatannya.
\end{enumerate}

Hasil pencarian jarak dan pemfaktoran menunjukkan beberapa faktor yang dominan:

\begin{center}
  \small
  \begin{tabular}{|c|c|}
    \hline
    Faktor & Frekuensi Muncul \\
    \hline
    2 & 2144 \\
    5 & 1898 \\
    10 & 1866 \\
    4 & 1014 \\
    20 & 888 \\
    \hline
  \end{tabular}
\end{center}

Tiga faktor teratas yang paling sering muncul berturut-turut adalah 2, 5, dan 10. Mengambil nilai maksimum dari sekumpulan faktor dengan frekuensi tertinggi ini memberikan tebakan panjang kunci yang sangat meyakinkan, yaitu \textbf{10}.

\subsubsection*{Tahap 2: \textit{Frequency Analysis}}

Setelah panjang kunci $L=10$ diketahui, \textit{ciphertext} dapat dibagi menjadi $L$ bagian (\textit{sub-text}), di mana \textit{sub-text} ke-$i$ berisi huruf-huruf pada posisi $i, i+L, i+2L, \dots$ dari \textit{ciphertext}. Karena setiap \textit{sub-text} dienkripsi dengan satu karakter kunci yang konstan, \textit{cipher} ini dapat dipandang sebagai gabungan $L$ buah sandi \textbf{Caesar \textit{Cipher}}.

\begin{enumerate}
  \item Untuk setiap \textit{sub-text}, hitung frekuensi kemunculan hurufnya.
  \item Lakukan uji coba pergeseran (\textit{shift}) dari 0 hingga 25.
  \item Untuk setiap pergeseran, hitung \textit{dot product} antara frekuensi huruf yang dihasilkan dengan probabilitas frekuensi huruf tunggal (\textit{monogram}) bahasa Inggris (\texttt{MONOGRAM} pada \texttt{src/frequencies.py}).
  \item Pilih pergeseran dengan nilai \textit{dot product} terbesar sebagai huruf ke-$i$ dari kunci (\texttt{VigenereSolver.guess\_key}).
\end{enumerate}

Berdasarkan analisis tersebut, kunci sepanjang 10 karakter berhasil ditebak, yaitu \textbf{\texttt{UNDERNIGHT}}.

\subsection*{Artefak}

Skrip \textit{solver} yang digunakan untuk melakukan \textit{Kasiski Examination}, menemukan kunci dengan analisis frekuensi, mendekripsi \textit{ciphertext}, serta menyimpannya ke berkas log.

\lstinputlisting[language=Python, basicstyle=\ttfamily\small, frame=single,
  title=\texttt{src/22/solver.py}]{../../src/22/solver.py}

\subsection*{\textit{Plaintext} Hasil Dekripsi}

\ciphertextfile{../../src/22/log/plain/2026-09-14_225345.txt}

\subsection*{\textit{Plaintext} dengan Format Asli}

\begin{tcolorbox}[breakable, colback=white, colframe=black!80,
  boxrule=0.6pt, arc=0mm, left=3mm, right=3mm, top=2mm, bottom=2mm]
\ttfamily\small\raggedright
Summer is drawing to a close. Soon, chirping crickets will take over for buzzing cicadas.\par
Gensokyo\textquotesingle{}s hot nights are cooling down, making it a comfortable time for humans and youkai alike.\par
Times are peaceful as usual.\par
From a human perspective, leastwise.\par
\medskip
\medskip
There is an old mansion located at the edge of Gensokyo. Its appearance radiates the feeling of its history and a kind of atmosphere that rejects visitors. Items that look as though they might be from the outer world have somehow found their way into this house. Strange machines, books that make no sense, magazines.\par
One such object, an iron box that was once a television, now serves the role of spiritual storage.\par
\textquotedbl{}Spirits can live well in something that used to display human shapes,\textquotedbl{} she\textquotesingle{}d say to her shikigami.\par
This is the home of Yukari Yakumo, the youkai of boundaries.\par
Ever since becoming aware of a slight abnormality in Gensokyo, she hasn\textquotesingle{}t been able to sleep well, even during the day.\par
The enemy was yet to be seen, but she was certain that it must have taken someone powerful to pull that off. She doesn\textquotesingle{}t get out often, however, and it will be very bothersome for her to take an action alone.\par
\textquotedbl{}Hey, I can rope her in and let her handle this problem.\textquotedbl{}\par
So she set out for the shrine, another building located at the edge of Gensokyo.\par
There\textquotesingle{}s a human that she knows there.\par
It\textquotesingle{}s a carefree person who always happens to be bored,\par
so she\textquotesingle{}ll gladly undertake the task of resolving whatever trouble there is.\par
\medskip
\medskip
It smells ominous. It is said that the forest eats humans. Humans usually maintain a safe distance from it. Its atmosphere constantly brims with an unearthly evil.\par
The Forest of Magic, a place where all the evil in Gensokyo naturally gathers.\par
Within that forest, small human figures are gathered into a small building.\par
Human figures smaller than humans.\par
The rainbow-{}colored puppeteer, Alice Margatroid, is reading books amidst a mountain of dolls.\par
\textquotedbl{}How can humans possibly not be aware of this huge disaster?\textquotedbl{}\par
\textquotedbl{}Dunno.\textquotedbl{}\par
If nothing is done, we won\textquotesingle{}t be able to enjoy that. She isn\textquotesingle{}t the type of person who gets out often, but no one else was interested in the problem, so she set out to investigate for herself.\par
Or rather, she tried to.\par
\textquotedbl{}I don\textquotesingle{}t like this, they\textquotesingle{}re used to this kind of thing, so they should do it.\textquotedbl{}\par
\textquotedbl{}Yeah, really.\textquotedbl{}\par
She knew nothing about the enemy and had no idea how to get started.\par
In the end, she decided to visit a human who also lived in the forest, to let her take care of the issue.\par
She brought a few books along......\par
Grimoires. Books that any human would find extremely rare and difficult to get ahold of.\par
That person would never be able to refuse such an offer.\par
\medskip
\medskip
\textquotedbl{}Sakuya, where are you?\textquotedbl{}\par
There is a European mansion located next to a lake; a scarlet mansion. A shrill voice was echoing all through its halls, as usual.\par
Between a lake of glistening white and a forest of verdant green stood a mansion of unyielding scarlet. It should be an awkward sight, but strangely, it\textquotesingle{}s very soothing.\par
It is said that time stops flowing within this mansion, the Scarlet Devil Mansion. It\textquotesingle{}s not a metaphor.\par
Remilia Scarlet, a vampire, was looking for her personal maid.\par
\textquotedbl{}Have you already done that as I asked?\textquotedbl{}\par
\textquotedbl{}I beg your pardon, my Lady, but I couldn\textquotesingle{}t get what actually you meant to...\textquotedbl{}\par
Unfortunately, she was apparently unsuccessful at communicating with the human before her.\par
\textquotedbl{}Fine! I\textquotesingle{}ll do it myself, Sakuya, so you do daily stu...\par
Well, whatever you want.\textquotedbl{}\par
Obviously, she had not been ordered to stay home. The maid has no choice but to attend her mistress as a protection charm. \textquotedbl{}You know that you won\textquotesingle{}t be able to do anything after sunrise,\textquotedbl{} she said in her heart.\par
It\textquotesingle{}s so peaceful and nothing in particular is happening, so she\textquotesingle{}ll just go back home once the mistress gets tired, the maid thought lightly. Of course, that too was left unspoken.\par
\medskip
\medskip
No other place in Gensokyo can match the silence of this place. However, it isn\textquotesingle{}t bleak. It has a peaceful atmosphere like a sleeping soul. There is no sign of any disturbance; nature\textquotesingle{}s bounty and refreshing wind is the sole source of sound.\par
This is the Netherworld, the land of the dead.\par
There isn\textquotesingle{}t a living human anywhere to be found. The ghosts there, however, enjoy a lively afterlife.\par
\textquotedbl{}So, Lady Yuyuko\textquotesingle{}s not aware of this?\textquotedbl{}\par
The most vivid and largest place in this quiet world. The Hakugyokurou.\par
Youmu Konpaku, a gardener, was pondering whether to tell her mistress about a most recent incident.\par
At that moment, she found her lady heading toward her. Good timing.\par
\textquotedbl{}Ah, Lady Yuyuko...\textquotedbl{}\par
\textquotedbl{}Youmu, are you still ignoring that?\textquotedbl{}\par
\textquotedbl{}Eh? ... What do you mean by \textquotesingle{}that\textquotesingle{}?\textquotedbl{}\par
\textquotedbl{}What, you\textquotesingle{}re not even aware of it?\par
That\textquotesingle{}s why you\textquotesingle{}re ragged about your dullness.\textquotedbl{}\par
She doesn\textquotesingle{}t remember being ragged, but it seems that her lady was aware of the trouble as well.\par
\textquotedbl{}So you were talking about the Moon? I am certainly aware of that.\par
You only said that, and so suddenly...\textquotedbl{}\par
\textquotedbl{}It looks like no one\textquotesingle{}s taking action, so why don\textquotesingle{}t you have a go, Youmu?\textquotedbl{}\par
\textquotedbl{}Ehh, why so?\textquotedbl{}\par
\textquotedbl{}Just kidding. You\textquotesingle{}re too unreliable anyway.\par
That human we saw the other day was better... well, I will do it then.\textquotedbl{}\par
\textquotedbl{}No way, please don\textquotesingle{}t be such a meanie... I shall go too.\textquotedbl{}\par
\textquotedbl{}I really meant what I said about unreliability, though.\textquotedbl{}\par
Yuyuko Saigyouji, the ghostly mistress of the Saigyouji family, ragged Youmu.\par
\textquotedbl{}Wait, do you have any idea where to go, my Lady?\textquotedbl{}\par
\textquotedbl{}Of course, a lot of them. As we shoot stuff down, we\textquotesingle{}ll eventually find a lead.\textquotedbl{}\par
\textquotedbl{}That\textquotesingle{}s where you\textquotesingle{}re wrong, Lady Yuyuko.\par
You rely too much on raw power and have bad aim, so it takes even longer.\par
You should concentrate your attack more, look...\textquotedbl{}\par
\textquotedbl{}Your back is so undefended, Youmu.\textquotedbl{}\par
Yuyuko was really worried about Youmu going out alone, so she decided to take action by herself.\par
Since the opponent is powerful enough to cause such a disaster, a team of two wouldn\textquotesingle{}t be a bad idea.\par
\medskip
\medskip
It was peaceful.\par
Or, it looked peaceful.\par
However, all youkai were upset.\par
What\textquotesingle{}s really happening is that, without anyone taking notice, suddenly...\par
The full moon has disappeared from Gensokyo\textquotesingle{}s night.\par
Even on the night of a supposed full moon, a bit of missing light prevents it from being full. It\textquotesingle{}s no wonder that humans aren\textquotesingle{}t aware of the situation, though; it\textquotesingle{}s only the tiniest bit of difference.\par
For non-{}humans, however, a moon that isn\textquotesingle{}t full doesn\textquotesingle{}t function at all.\par
It\textquotesingle{}s an especially critical problem for those who dislike sunlight.\par
A human and youkai, a team of two, ventured out into Gensokyo\textquotesingle{}s midnight hours\par
to retrieve the missing fragment of the moon and restore the full moon.\par
They will stop the night if need be,\par
even if it becomes an imperishable night.\par
------ As summer draws to a close, not much time remains before the harvest moon.\par
The team of human and youkai stop the night.\par
\medskip
\medskip
A glowing, ancient, virgulate object.\par
A pile of round objects in the sight.\par
A small gem. A shining orb. A fragile soul. And the largest sphere.\par
\textquotedbl{}How is she doing these days?\textquotedbl{} she wondered as she stared at these objects.\par
Time does not flow in this place. History repeats itself.\par
She, too, was in Gensokyo.\par
\bigskip
Sumber:\par
\hspace*{1.5em}- \url{https://en.touhouwiki.net/wiki/Imperishable_Night/Story/Prologue}
\end{tcolorbox}


\subsection*{Kunci}

Panjang kunci: \textbf{10}

Kunci hasil rekonstruksi: \textbf{\texttt{UNDERNIGHT}}

\subsection*{Referensi}

Tabel frekuensi \textit{monogram} bahasa Inggris yang digunakan pada Tahap 2 mengacu pada \cite{cornell-freq}. Referensi mengenai konsep dasar dekripsi Vigen\`ere \textit{Cipher} mengacu pada \cite{vigenere-gfg}. Konsep Vigen\`ere \textit{Cipher} dan metode Kasiski mengacu pada bahan kuliah \cite{munir-klasik2}.
